% ===== 原文 pp.92–95（PDF p.102–105） =====
% [原文 p.92]
的修正 A-模型作用量。按构造，$S_{top}$ 被 BRST 变分 $\delta_Q$ 湮灭；快速检查一下上一节的论证即可看出，它也在映射 $\Phi$ 与联络 $\omega$ 的任意无穷小变分下不变。在这个意义上，它代表作用量中的一个拓扑项，类似于 (7.3) 中的拓扑项。

我们可以把这个“等变”拓扑项解释为：它测度的是靶空间上某个等变上同调类的拉回。设 $\eta\in\Omega^2_{U(1)}(X)$ 是由
\[
\eta=\Omega-H\zeta
\]
给出的等变 2-形式，其中 $\zeta$ 是 $S(\fg^*)$ 的生成元，$\fg=\fu(1)$。$U(1)$ 在 $X$ 上的作用是哈密顿作用、矩映射为 $H$ 这一假定，蕴含 $\eta$ 在作用在等变形式上的 Cartan 微分 $\dd_C$ 下闭：
\[
\dd_C\eta=(\dd-\iota_V\zeta)(\Omega-H\zeta)=-(\dd H)\zeta-(\iota_V\Omega)\zeta=0 .
\]
如附录 C 所述，可以用截面 $\Phi:\Sigma\to E$ 与联络 $\omega$ 把 $X$ 上的等变 2-形式 $\eta$ 拉回为 $\Sigma$ 上的普通 2-形式 $\tilde\eta$，再在 $\Sigma$ 上积分得到一个拓扑项。

确切地说，我们先利用协变导数 $\dd_\omega\varphi$ 把 $\eta$ 拉回为 $P$ 上的普通 2-形式 $\eta(\varphi,\omega)$，其中 $\varphi:P\to X$ 是截面 $\Phi$ 对应的 $G$-等变映射。形式 $\eta(\varphi,\omega)$ 是 basic 的，因此它下降到 $\Sigma$ 上一个整体定义、闭的 2-形式 $\tilde\eta(\varphi,\omega)$，在每个邻域 $U^{(\alpha)}$ 上由 $\tilde\eta(\varphi,\omega)=s^{(\alpha)*}\eta(\varphi,\omega)$ 给出。由附录 C 的讨论可以保证，该形式的积分给出一个数
\[
\int_\Sigma\tilde\eta(\varphi,\omega),
\]
它只依赖于 $\Phi$ 与 $\omega$ 的同伦类以及 $\eta$ 的等变上同调类。在局部邻域 $U^{(\alpha,\gamma)}$ 上，我们可以把 $\tilde\eta$ 写成
\begin{align*}
\tilde\eta(\varphi,\omega)&=s^{(\alpha)*}\eta(\varphi,\omega)\\
&=s^{(\alpha)*}\Big[(\dd_\omega\varphi)^*\Omega_X-(H\circ\varphi)F_\omega\Big]\\
&=\big[\dd_\omega\varphi\circ\dd s^{(\alpha)}\big]^*\Omega_X-\big(H\circ\varphi\circ s^{(\alpha)}\big)s^{(\alpha)*}F_\omega\\
&=\big[\dd_\omega\Phi^{(\alpha)}\big]^*\Big(\frac12\Omega^{(\gamma)}_{IJ}\dd\phi^{(\gamma)I}\wedge\dd\phi^{(\gamma)J}\Big)-\big(H\circ\Phi^{(\alpha)}\big)\dd\omega^{(\alpha)}\\
&=\frac12\big(\Omega^{(\gamma)}_{IJ}\circ\Phi^{(\alpha)}\big)\mathcal{D}\phi^I_{(\alpha,\gamma)}\wedge\mathcal{D}\phi^J_{(\alpha,\gamma)}-\big(H\circ\Phi^{(\alpha)}\big)\dd\omega^{(\alpha)} .
\end{align*}
我们认出，上式的积分正是前面写下的 $S_{top}$。

% [Tutor Note: p.92 “$S_{top}$ 对 $\Phi$ 与 $\omega$ 的任意变分不变”这一步是整节解释的关键，补出论证：把 $\tilde\eta$ 的表达式取变分，$\delta\Phi$ 与 $\delta\omega$ 只会使 $\tilde\eta$ 改变一个全微分（因为 $\tilde\eta$ 的 $\dd$ 闭性来自 $\dd_C\eta=0$ 中 $\zeta\to\dd\omega$ 的替换），在闭曲面 $\Sigma$ 上积分为零，故 $S_{top}$ 只依赖同伦类。这也解释了正文开头“by construction”之外的“a quick inspection”究竟在查什么。]
% [Tutor Note: p.92 与 p.94 用了同一个符号 $\eta$ 表示两个不同的对象：p.92 的 $\eta=\Omega-H\zeta$ 只是任意的等变 2-形式（用来定义 $S_{top}$ 的等变类），p.94 的 $\eta=\frac{1}{2\pi}\Omega-\frac{1}{4\pi}\zeta\Delta H$ 则是专门的\emph{等变第一陈类}。两处不可混用；本译文中 $\tilde\eta$ 一律指拉回到 $\Sigma$ 上的普通 2-形式（原文的 $\tilde\eta$）。另注意 p.92 的 $\dd_C\eta=0$ 只有在 (7.5) 取负号约定（$\dd H=-\iota_V\Omega$）时才成立，见 p.88 批注。]
% [Tutor Note: p.92 那条五行的链式推导可逐行核对：第二行到第三行用了 $F_\omega=\dd\omega$，且 $\varphi\circ s^{(\alpha)}=\Phi^{(\alpha)}$；第三行到第四行把 $\Omega_X$ 展开为 p.86 约定的 $\frac12\Omega_{IJ}\dd\phi^I\wedge\dd\phi^J$，并把 $\dd s^{(\alpha)}$ 并入协变导数得 $\mathcal{D}\phi^I_{(\alpha,\gamma)}$；最后一行与 §7.3 开头 $S_{top}$ 的被积式逐项相同（相差整体的积分号）。若核对时发现 $s^{(\alpha)*}F_\omega$ 与 $\dd\omega^{(\alpha)}$ 的标签 $(\alpha)$、$(\gamma)$ 对不上，注意：$\omega^{(\alpha)}$ 是同一个整体联络 $\omega$ 在 chart $U^{(\alpha)}$ 上的表达式，而 $\Omega^{(\gamma)}$ 是靶空间 chart $U^{(\gamma)}$ 上的坐标表达式。]

% [原文 p.93]
\section{局部可观测量}

修正 A-模型的局部可观测量谱相当直截了当：它就是把 de Rham 上同调限制到其 $U(1)$ 不变子复形上。

更确切地说，考虑靶空间上闭的微分 $k$-形式 $\rho$ 构成的空间，其中 $\rho$ 是 $U(1)$ 不变的，即 $\mathcal{L}_V\rho=0$。利用李导数的 Cartan 恒等式立即看出，外微分 $\dd\rho$ 也是 $U(1)$ 不变的。因此我们可以一致地把 $X$ 上的 de Rham 复形限制到其 $U(1)$ 不变子复形。这个子复形的上同调称为 $X$ 的 $U(1)$ 不变上同调，记作 $(H^*X)^{U(1)}$。

% [Tutor Note: p.93 的 $(H^*X)^{U(1)}$ 是 $\Omega^*(X)$ 的 $U(1)$ 不变子复形的上同调，\emph{不是}等变上同调 $H^*_{U(1)}(X)$（后者是 p.92 的 Cartan 复形 $(\Omega^*(X)\otimes S(\fg^*))^{U(1)}$ 的上同调）。两者关系可用例子核对：若 $U(1)$ 作用在 $X$ 上自由，则 $H^*_{U(1)}(X)\cong H^*(X/U(1))$，而 $(H^*X)^{U(1)}\subset H^*(X)$ 只是不变部分；对 $X=\C P^1$ 上的标准旋转，$H^0$ 与 $H^2$ 都不动，于是 $(H^*X)^{U(1)}=H^*(X)$，可作为最简单的检验。]

\para{鬼数反常}

现在计算 $U(1)$ 鬼数对称性的量子力学反常；它的消失将对上述可观测量中哪些能有非零期望值施加选择规则。反常来自路径积分测度中鬼数为 1 的费米子（$\chi$ 与 $\rho$）的零模数目与鬼数为 $-1$ 的费米子（$\tilde\chi$ 与 $\tilde\rho$）的零模数目不匹配。（注意：此处符号 $\chi$ 指我们先前称为 $\chi$ 的场的 $(1,0)$ 部分，而符号 $\tilde\chi$ 指它的 $(0,1)$ 部分。）

零模的不匹配对应于费米子动能项中出现的椭圆算子 $\mathcal{D}$ 的指标（的两倍）：
\begin{align*}
k&=2\dim\ker\Big\{\mathcal{D}_{\bar z}:\Omega^0\big(\Sigma,\Phi^*V^{1,0}E\big)\to\Omega^{0,1}\big(\Sigma,\Phi^*V^{1,0}E\big)\Big\}\\
&\quad-2\dim\ker\Big\{\mathcal{D}_z:\Omega^0\big(\Sigma,\Phi^*V^{0,1}E\big)\to\Omega^{1,0}\big(\Sigma,\Phi^*V^{0,1}E\big)\Big\}\\
&=2\dim\ker\mathcal{D}_{\bar z}-2\dim\ker\mathcal{D}_{\bar z}^\dagger\\
&=2\operatorname{ind}\mathcal{D}_{\bar z}.
\end{align*}
由 Hirzebruch-Riemann-Roch 定理，这个作用在全纯向量丛 $\Phi^*V^{1,0}E$ 上的 Cauchy-Riemann 算子的指标可由其第一陈数给出：
\[
\operatorname{ind}\mathcal{D}_{\bar z}=n(1-g_\Sigma)+\int_\Sigma c_1\big(\Phi^*V^{1,0}E\big).
\]
第一陈类还可以进一步 [12] 解释为切丛 $TX$ 上等变第一陈类在截面 $\Phi$ 下的拉回：
\[
c_1\big(\Phi^*V^{1,0}E\big)=\tilde\eta(\varphi,\omega)
\]

% [Tutor Note: p.93 指标公式里的因子 $2$ 对应原文的“twice the index”：第二行里的 $\ker\mathcal{D}_z$（作用在取值于 $V^{0,1}E$ 的 0-形式、映到 $\Omega^{1,0}$）经 Hodge 星把 $(1,0)$-形式对偶为 $(0,1)$-形式后，其维数等于伴随算子 $\mathcal{D}_{\bar z}^\dagger$ 的核维数；而 $\operatorname{ind}\mathcal{D}_{\bar z}=\dim\ker\mathcal{D}_{\bar z}-\dim\ker\mathcal{D}_{\bar z}^\dagger$，故 $k=2\operatorname{ind}\mathcal{D}_{\bar z}$。用靶空间为单点（$n=0$）检验：$k=0$，反常消失，与“点靶空间的 A-模型只有常数可观测量”相符。]
% [Tutor Note: p.93 末 $c_1(\Phi^*V^{1,0}E)=\tilde\eta(\varphi,\omega)$ 中的 $\tilde\eta$ 是 p.94 那个等变第一陈类的拉回（不是 p.92 的一般 $\eta$ 的拉回）。这条等式与附录 C 的“沿纤维积分/Cartan 模型”构造配套使用；它的一个可直接检验的后果在 p.99：当 $X$ 是等变 Calabi-Yau 时，$\int_\Sigma\tilde\eta$ 与 $(F,F_{\mathbb R})$ 的 Maslov 指标相等，两处必须用同一个 $\tilde\eta$。]

% [原文 p.94]
其中
\[
\eta=\frac{1}{2\pi}\Omega-\frac{1}{4\pi}\zeta\,\Delta H\quad\in\Omega^2_{U(1)}(X)
\]
是等变第一陈类，而 $\tilde\eta(\varphi,\omega)$ 是它的拉回（记号同上一节）。这里 $\Delta H$ 是 Hamilton 函数的拉普拉斯算子，$\Omega$ 是 $X$ 上的 Kähler 形式。

% [Tutor Note: p.94 的 $\eta=\frac{1}{2\pi}\Omega-\frac{1}{4\pi}\zeta\Delta H$ 与 p.92 的 $\eta=\Omega-H\zeta$ 形状相同但系数不同：前者是 $\mathrm{ch}$ 型的归一（$1/2\pi$ 来自陈类的标准归一，$\Delta H$ 来自等变陈类的曲率修正），后者是任意等变形式。核对方法：在 $X=\C P^1$、$U(1)$ 为绕 $S^2$ 轴旋转、$H$ 为高度函数的情形下，$\Delta H$ 正比于 $H$，把 $\eta$ 在 $\Sigma$ 上积分可复现第一陈数的整数归一。]

\section{边界条件}

现在考虑把修正 A-模型放在带非平凡边界 $\partial\Sigma$ 的黎曼面 $\Sigma$ 上，边界由 $h_\Sigma$ 个微分同胚于圆周的分支组成，它们从 $\Sigma$ 的定向自然继承定向。正如我们将在第 9 章进一步讨论的，必须沿 $\partial\Sigma$ 对场施加局部边界条件，以确保作用量变分中的边界项对所有允许变分恒为零。此外，要求边界条件保持体理论的对称性 $\delta_Q$ 会对边界条件施加额外约束；下面我们详细看这一点。

\noindent\emph{$U(1)$ 不变边界条件}

我们写出一组充分的局部边界条件（暂时搁置必要性的问题）。我们写下的条件隐含地理解为沿 $\partial\Sigma$ 施加。

我们可以使用主丛 $P$ 在 $\Sigma$ 上的任意一组相容平凡化 chart；为简单起见，设 $\partial\Sigma$ 的每个分支恰好落在一个 chart $U^{(\alpha)}$ 内。与普通 A-模型一样，我们施加这样的条件：$\sigma$-模型映射 $\Phi^{(\alpha)}$ 取值于子流形 $Y^{(\alpha)}\subset X$，其中 $Y^{(\alpha)}$ 关于辛形式 $\Omega$ 是拉格朗日的。我们还进一步要求 $Y^{(\alpha)}$ 是 $U(1)$ 不变的，意思是生成性的 Killing 向量场 $V$ 在 $Y^{(\alpha)}$ 的所有点处切于 $Y^{(\alpha)}$。由于 $Y^{(\alpha)}$ 是迷向的，$U(1)$ 不变性蕴含 Hamilton 函数 $H$ 在该子流形上局部为常数。（这是因为 $H$ 沿切向量 $W\in TY$ 的变化率由 $W^I\partial_IH=\Omega_{IJ}V^IW^J=0$ 给出。）反之，$Y^{(\alpha)}$ 是余迷向的这一事实意味着 $H$ 为常数蕴含 $U(1)$ 不变性。因此，$U(1)$ 不变性的假定等价于假定 $H$ 在边界 $\partial\Sigma$ 的每个分支上取常值。对这样的子流形，就不再需要标号 $(\alpha)$。

$\sigma$-模型映射的允许变分 $\delta\Phi$ 是切于该子流形的变分，即
\[
\delta\phi^I\frac{\partial}{\partial\phi^I}\in\Omega^0(\partial\Sigma,\Phi^*TY)\subset\Omega^0(\partial\Sigma,\Phi^*TX).
\]
预见到 BRST 不变性，我们希望 $\delta=\delta_Q$ 是允许变分，因此施加

% [Tutor Note: p.94 的“$Y$ 拉格朗日且 $U(1)$ 不变 $\Leftrightarrow$ $H|_Y$ 局部常值”要把两个方向都补出来：$Y$ 迷向给出 $\Omega|_{TY}=0$，结合 $V$ 切于 $Y$ 得 $\iota_V\Omega|_{TY}=0$；而 $\dd H|_Y=-(\iota_V\Omega)|_Y$ 在 $TY$ 上为零，故 $H|_Y$ 局部常值。反向：$H$ 在 $Y$ 上常值时 $\iota_V\Omega|_Y=0$，即 $V$ 在 $TY$ 的辛正交补里；$Y$ 拉格朗日意味着 $(TY)^\perp=TY$，于是 $V$ 切于 $Y$。原文 p.94 该行有笔误“the the symplectic form”，重复了冠词，不影响数学。]
% [Tutor Note: p.94 这里 A-膜（brane）的术语要分清：普通 A-模型允许的边界条件是拉格朗日子流形（Lagrangian），而在扭转过（带 $B$-场/等变拓扑项）的理论里更一般的对象是\emph{余迷向}（coisotropic）子流形；本段同时用到“迷向”与“余迷向”两个性质，是因为 $Y$ 被假定为拉格朗日（既迷向又余迷向），故两个方向的蕴含都成立。精确地说：迷向性给出“$U(1)$ 不变 $\Rightarrow$ $H$ 在 $Y$ 上局部常值”，余迷向性给出反方向；若只保留余迷向而丢掉迷向，则后一方向不再成立。]

% [原文 p.95]
费米场 $\chi$ 也切于该子流形的条件。

对体作用量作场的无穷小变分并代入体运动方程，剩下的边界项为
\[
\delta S=\int_{\partial\Sigma}\star\,g_{IJ}\dd\sigma^\mu n_\mu\Big\{\mathcal{D}_\mu\phi^I\delta\phi^J+2\ii\rho^I_\mu\delta^{(c)}\chi^J-2\ii\tilde\rho^I_\mu\delta^{(c)}\chi^J\Big\},
\]
其中 $\star$ 是 $\Sigma$ 上的 Hodge 星，而
\[
\delta^{(c)}\chi^J\equiv\delta\chi^J+\Gamma^J{}_{KL}\delta\phi^K\chi^L
\]
是费米子的\emph{协变变分}。与体运动方程相容要求该项对所有允许变分恒为零。为保证这一点，我们施加协变 Neumann 条件：$n^\mu\mathcal{D}_\mu\phi^I$ 取值于法子丛 $(TY)^\perp$（关于度规 $g$），其中 $n$ 是边界的法向（关于度规 $h$）。同样，我们要求 $n^\mu(\rho_\mu-\tilde\rho_\mu)$ 垂直于该子流形。总之，我们施加如下边界条件：
\begin{equation}
\begin{aligned}
\Phi&\in Y\subset X\\
\iota_n(\dd_\omega\Phi)&\in\Omega^0\big(\partial\Sigma,\Phi^*(TY)^\perp\big)\\
\chi&\in\Omega^0\big(\partial\Sigma,\Phi^*(TY)\big)\\
\iota_n(\rho-\tilde\rho)&\in\Omega^0\big(\partial\Sigma,\Phi^*(TY)^\perp\big).
\end{aligned} \tag{7.9}
\end{equation}
在离壳处理时，只需把上面第二行中的 $\dd_\omega\Phi$ 换成 $P-\tilde P$。

现在验证这些边界条件 BRST 不变。在离壳情形利用 $\delta_Q^2=0$，可得作用量变分中的边界项
\[
\delta_QS=\delta_QS_{top}=-\int_{\partial\Sigma}\Omega_{IJ}\mathcal{D}_\mu\phi^I\chi^J,
\]
它之所以消失，是因为 $V$、$\chi$ 与 $t^\mu\partial_\mu\phi^I$ 都切于该子流形（其中 $t$ 切于边界）。

最后，考虑把世界面度规变化 $h_{\mu\nu}\to h_{\mu\nu}+\delta h_{\mu\nu}$ 的效果。为了继续得到一个拓扑场论，作用量在 $\delta_h$ 下必须只改变一个 BRST-正合量。快速检查离壳变分 (7.2) 即可看出，度规 $h$ 在其中并不出现，因此 $\delta_h$ 与 BRST 变分 $\delta_Q$ 对易。于是
\[
\delta_hS=\delta_Q\delta_hV+\delta_hS_{top}=(\text{Q-exact})-\int_{\partial\Sigma}H(\Phi)\delta_h\omega
\]

% [Tutor Note: p.95 (7.9) 的四行条件可逐条对照边界项的消失来核对：第 1 行（$\Phi$ 取值于 $Y$）使 $\delta\phi^J\in TY$；第 2 行（法向协变导数落在 $(TY)^\perp$）使 $\mathcal{D}_\mu\phi^I\delta\phi^J$ 项因 $\star n^\mu\mathcal{D}_\mu\phi$ 与 $T Y$ 正交而消失；第 3、4 行（$\chi$ 切向、$\iota_n(\rho-\tilde\rho)$ 法向）使后两项相消。注意第 2 行写的是 $\iota_n(\dd_\omega\Phi)$ 而不是 $n^\mu\mathcal{D}_\mu\phi^I$：两者在 $\partial\Sigma$ 上对同一对象的两种写法，离壳时需要用 $P-\tilde P$（$P$ 与 $\tilde P$ 分别是 $\mathcal{D}$ 的 $(0,1)$、$(1,0)$ 部分）。]
% [Tutor Note: p.95 末 $\delta_hS$ 的正合性要求把 $-\int_{\partial\Sigma}H(\Phi)\delta_h\omega$ 抵消掉，这就是下一页引入 $S_{\partial\Sigma}=\int_{\partial\Sigma}H(\Phi)\omega$ 的动机。核对符号：$\delta_hS_{\partial\Sigma}=+\int_{\partial\Sigma}H\delta_h\omega$，与上式恰好相消；此外 $\delta_QS_{\partial\Sigma}=\int_{\partial\Sigma}(\delta_QH)\omega=0$（因 $\chi$ 切于 $Y$ 且 $\dd H|_Y=0$），故不破坏 BRST 不变性。]
